\begin{titlepage}
\thispagestyle{empty}
\vspace*{3.2cm}
{\Huge\bfseries Cartly}\par
\vspace{6pt}
{\LARGE An Engineering Reconstruction of a\\Marketplace-Scale Shopping Experience}\par
\vspace{14pt}
\rule{\linewidth}{0.6pt}\par
\vspace{6pt}
{\large Amazon.com Clone --- 24-Hour Software Engineering Challenge}\par
{\large Engineering Report}\par
\vspace{2.2cm}

\begin{tabular}{@{}p{3.4cm}p{10.5cm}@{}}
\textbf{Author} & Nitish R.G. \\[3pt]
\textbf{Date} & 3 October 2026 \\[3pt]
\textbf{Repository} & \code{github.com/NITISH-R-G/amazon-clone} (branch \code{main}) \\[3pt]
\textbf{Deployment} & \code{amazon-clone-six-sigma-27.vercel.app} (Vercel, HTTPS) \\[3pt]
\textbf{Stack} & Next.js 16, React 19, TypeScript (strict), PostgreSQL with Drizzle ORM, Stripe (test mode), Vitest, Playwright \\[3pt]
\textbf{Scope} & A fictional marketplace: invented brands, sellers, products and reviews. No Amazon data or assets are used. \\
\end{tabular}

\vfill
{\small\textbf{Reading guide.} Throughout this report, claims are labelled by their evidence:
\tagimpl{} (present in the source), \tagtest{} (covered by an automated test), \tagplan{} (considered or designed only) and
\tagunver{} (could not be confirmed from the repository).}
\end{titlepage}
